\documentclass[a4paper,12pt,twoside]{memoir}

% Castellano
\usepackage[spanish,es-tabla]{babel}
\selectlanguage{spanish}
\usepackage[utf8]{inputenc}
\usepackage[T1]{fontenc}
\usepackage{lmodern} % Scalable font
\usepackage{microtype}
\usepackage{placeins}

\RequirePackage{booktabs}
\RequirePackage[table]{xcolor}
\RequirePackage{xtab}
\RequirePackage{multirow}

% Links
\usepackage[colorlinks]{hyperref}
\hypersetup{
	allcolors = {red}
}

% Ecuaciones
\usepackage{amsmath}

% Rutas de fichero / paquete
\newcommand{\ruta}[1]{{\sffamily #1}}

% Párrafos
\nonzeroparskip


% Imagenes
\usepackage{graphicx}
\newcommand{\imagen}[2]{
	\begin{figure}[!h]
		\centering
		\includegraphics[width=0.9\textwidth]{#1}
		\caption{#2}\label{fig:#1}
	\end{figure}
	\FloatBarrier
}
\newcommand{\imagencite}[3]{
	\begin{figure}[!h]
		\centering
		\includegraphics[width=0.9\textwidth]{#1}
		\caption[#2]{#2 \cite{#3}}\label{fig:#1}
	\end{figure}
	\FloatBarrier
}

\newcommand{\imagenflotante}[2]{
	\begin{figure}%[!h]
		\centering
		\includegraphics[width=0.9\textwidth]{#1}
		\caption{#2}\label{fig:#1}
	\end{figure}
}

% Pie de página con enlace
\newcommand{\foothref}[2]{
	\href{#2}{#1}\footnote{#1: \href{#2}{#2}}
}



% El comando \figura nos permite insertar figuras comodamente, y utilizando
% siempre el mismo formato. Los parametros son:
% 1 -> Porcentaje del ancho de página que ocupará la figura (de 0 a 1)
% 2 --> Fichero de la imagen
% 3 --> Texto a pie de imagen
% 4 --> Etiqueta (label) para referencias
% 5 --> Opciones que queramos pasarle al \includegraphics
% 6 --> Opciones de posicionamiento a pasarle a \begin{figure}
\newcommand{\figuraConPosicion}[6]{%
  \setlength{\anchoFloat}{#1\textwidth}%
  \addtolength{\anchoFloat}{-4\fboxsep}%
  \setlength{\anchoFigura}{\anchoFloat}%
  \begin{figure}[#6]
    \begin{center}%
      \Ovalbox{%
        \begin{minipage}{\anchoFloat}%
          \begin{center}%
            \includegraphics[width=\anchoFigura,#5]{#2}%
            \caption{#3}%
            \label{#4}%
          \end{center}%
        \end{minipage}
      }%
    \end{center}%
  \end{figure}%
}

%
% Comando para incluir imágenes en formato apaisado (sin marco).
\newcommand{\figuraApaisadaSinMarco}[5]{%
  \begin{figure}%
    \begin{center}%
    \includegraphics[angle=90,height=#1\textheight,#5]{#2}%
    \caption{#3}%
    \label{#4}%
    \end{center}%
  \end{figure}%
}
% Para las tablas
\newcommand{\otoprule}{\midrule [\heavyrulewidth]}
%
% Nuevo comando para tablas pequeñas (menos de una página).
\newcommand{\tablaSmall}[5]{%
 \begin{table}
  \begin{center}
   \rowcolors {2}{gray!35}{}
   \begin{tabular}{#2}
    \toprule
    #4
    \otoprule
    #5
    \bottomrule
   \end{tabular}
   \caption{#1}
   \label{tabla:#3}
  \end{center}
 \end{table}
}

%
% Nuevo comando para tablas pequeñas (menos de una página).
\newcommand{\tablaSmallSinColores}[5]{%
 \begin{table}[H]
  \begin{center}
   \begin{tabular}{#2}
    \toprule
    #4
    \otoprule
    #5
    \bottomrule
   \end{tabular}
   \caption{#1}
   \label{tabla:#3}
  \end{center}
 \end{table}
}

\newcommand{\tablaApaisadaSmall}[5]{%
\begin{landscape}
  \begin{table}
   \begin{center}
    \rowcolors {2}{gray!35}{}
    \begin{tabular}{#2}
     \toprule
     #4
     \otoprule
     #5
     \bottomrule
    \end{tabular}
    \caption{#1}
    \label{tabla:#3}
   \end{center}
  \end{table}
\end{landscape}
}

%
% Nuevo comando para tablas grandes con cabecera y filas alternas coloreadas en gris.
\newcommand{\tabla}[6]{%
  \begin{center}
    \tablefirsthead{
      \toprule
      #5
      \otoprule
    }
    \tablehead{
      \multicolumn{#3}{l}{\small\sl continúa desde la página anterior}\\
      \toprule
      #5
      \otoprule
    }
    \tabletail{
      \hline
      \multicolumn{#3}{r}{\small\sl continúa en la página siguiente}\\
    }
    \tablelasttail{
      \hline
    }
    \bottomcaption{#1}
    \rowcolors {2}{gray!35}{}
    \begin{xtabular}{#2}
      #6
      \bottomrule
    \end{xtabular}
    \label{tabla:#4}
  \end{center}
}

%
% Nuevo comando para tablas grandes con cabecera.
\newcommand{\tablaSinColores}[6]{%
  \begin{center}
    \tablefirsthead{
      \toprule
      #5
      \otoprule
    }
    \tablehead{
      \multicolumn{#3}{l}{\small\sl continúa desde la página anterior}\\
      \toprule
      #5
      \otoprule
    }
    \tabletail{
      \hline
      \multicolumn{#3}{r}{\small\sl continúa en la página siguiente}\\
    }
    \tablelasttail{
      \hline
    }
    \bottomcaption{#1}
    \begin{xtabular}{#2}
      #6
      \bottomrule
    \end{xtabular}
    \label{tabla:#4}
  \end{center}
}

%
% Nuevo comando para tablas grandes sin cabecera.
\newcommand{\tablaSinCabecera}[5]{%
  \begin{center}
    \tablefirsthead{
      \toprule
    }
    \tablehead{
      \multicolumn{#3}{l}{\small\sl continúa desde la página anterior}\\
      \hline
    }
    \tabletail{
      \hline
      \multicolumn{#3}{r}{\small\sl continúa en la página siguiente}\\
    }
    \tablelasttail{
      \hline
    }
    \bottomcaption{#1}
  \begin{xtabular}{#2}
    #5
   \bottomrule
  \end{xtabular}
  \label{tabla:#4}
  \end{center}
}



\definecolor{cgoLight}{HTML}{EEEEEE}
\definecolor{cgoExtralight}{HTML}{FFFFFF}

%
% Nuevo comando para tablas grandes sin cabecera.
\newcommand{\tablaSinCabeceraConBandas}[5]{%
  \begin{center}
    \tablefirsthead{
      \toprule
    }
    \tablehead{
      \multicolumn{#3}{l}{\small\sl continúa desde la página anterior}\\
      \hline
    }
    \tabletail{
      \hline
      \multicolumn{#3}{r}{\small\sl continúa en la página siguiente}\\
    }
    \tablelasttail{
      \hline
    }
    \bottomcaption{#1}
    \rowcolors[]{1}{cgoExtralight}{cgoLight}

  \begin{xtabular}{#2}
    #5
   \bottomrule
  \end{xtabular}
  \label{tabla:#4}
  \end{center}
}


















\graphicspath{ {./img/} }

% Capítulos
\chapterstyle{bianchi}
\newcommand{\capitulo}[2]{
	\setcounter{chapter}{#1}
	\setcounter{section}{0}
	\chapter*{#2}
	\addcontentsline{toc}{chapter}{#2}
	\markboth{#2}{#2}
}

% Apéndices
\renewcommand{\appendixname}{Apéndice}
\renewcommand*\cftappendixname{\appendixname}

\newcommand{\apendice}[1]{
	%\renewcommand{\thechapter}{A}
	\chapter{#1}
}

\renewcommand*\cftappendixname{\appendixname\ }

% Formato de portada
\makeatletter
\usepackage{xcolor}
\newcommand{\tutor}[1]{\def\@tutor{#1}}
\newcommand{\course}[1]{\def\@course{#1}}
\definecolor{cpardoBox}{HTML}{E6E6FF}
\def\maketitle{
  \null
  \thispagestyle{empty}
  % Cabecera ----------------
\noindent\includegraphics[width=\textwidth]{cabecera}\vspace{1cm}%
  \vfill
  % Título proyecto y escudo informática ----------------
  \colorbox{cpardoBox}{%
    \begin{minipage}{.8\textwidth}
      \vspace{.5cm}\Large
      \begin{center}
      \textbf{TFG del Grado en Ingeniería Informática}\vspace{.6cm}\\
      \textbf{\LARGE\@title{}}
      \end{center}
      \vspace{.2cm}
    \end{minipage}

  }%
  \hfill\begin{minipage}{.20\textwidth}
    \includegraphics[width=\textwidth]{escudoInfor}
  \end{minipage}
  \vfill
  % Datos de alumno, curso y tutores ------------------
  \begin{center}%
  {%
    \noindent\LARGE
    Presentado por \@author{}\\ 
    en Universidad de Burgos --- \@date{}\\
    Tutores: \@tutor{}\\
  }%
  \end{center}%
  \null
  \cleardoublepage
  }
\makeatother

\newcommand{\nombre}{Iván Arjona Alonso} %%% cambio de comando

% Datos de portada
\title{Aplicación Web para la recopilación, tratamiento y visualización de datos públicos}
\author{\nombre}
\tutor{Dr. José Francisco Díez Pastor\\y Dr. Jesús Manuel Maudes Raedo}
\date{\today}

\begin{document}

\maketitle


\newpage\null\thispagestyle{empty}\newpage


%%%%%%%%%%%%%%%%%%%%%%%%%%%%%%%%%%%%%%%%%%%%%%%%%%%%%%%%%%%%%%%%%%%%%%%%%%%%%%%%%%%%%%%%
\thispagestyle{empty}


\noindent\includegraphics[width=\textwidth]{cabecera}\vspace{1cm}

\noindent D. José Francisco Díez Pastor y D. Jesús Manuel Maudes Raedo, profesores del departamento de Ingeniería Civil, área de Lenguajes y Sistemas Informáticos.

\noindent Expone:

\noindent Que el alumno D. \nombre, con DNI 71352655P, ha realizado el Trabajo final de Grado en Ingeniería Informática titulado ``Aplicación Web para la recopilación, tratamiento y visualización de datos públicos''. 

\noindent Y que dicho trabajo ha sido realizado por el alumno bajo la dirección del que suscribe, en virtud de lo cual se autoriza su presentación y defensa.

\begin{center} %\large
En Burgos, {\large \today}
\end{center}

\vfill\vfill\vfill

% Author and supervisor
\begin{minipage}{0.45\textwidth}
\begin{flushleft} %\large
Vº. Bº. del Tutor:\\[2cm]
D. José Francisco Díez Pastor
\end{flushleft}
\end{minipage}
\hfill
\begin{minipage}{0.45\textwidth}
\begin{flushleft} %\large
Vº. Bº. del Tutor:\\[2cm]
D. Jesús Manuel Maudes Raedo
\end{flushleft}
\end{minipage}
\hfill

\vfill

% para casos con solo un tutor comentar lo anterior
% y descomentar lo siguiente
%Vº. Bº. del Tutor:\\[2cm]
%D. nombre tutor


\newpage\null\thispagestyle{empty}\newpage




\frontmatter

% Abstract en castellano
\renewcommand*\abstractname{Resumen}
\begin{abstract}
En este primer apartado se hace una \textbf{breve} presentación del tema que se aborda en el proyecto.
\end{abstract}

\renewcommand*\abstractname{Descriptores}
\begin{abstract}
Palabras separadas por comas que identifiquen el contenido del proyecto Ej: servidor web, buscador de vuelos, android \ldots
\end{abstract}

\clearpage

% Abstract en inglés
\renewcommand*\abstractname{Abstract}
\begin{abstract}
A \textbf{brief} presentation of the topic addressed in the project.
\end{abstract}

\renewcommand*\abstractname{Keywords}
\begin{abstract}
keywords separated by commas.
\end{abstract}

\clearpage

% Indices
\tableofcontents

\clearpage

\listoffigures

\clearpage

\listoftables
\clearpage

\mainmatter
\capitulo{1}{Introducción}

La \textit{minería de datos} es una metodología para descubrir relaciones ocultas en grandes cantidades de datos utilizando técnicas de inteligencia artificial. Actualmente se está extendiendo bastante el uso de técnicas de minería de datos para realizar análisis sociológicos, como por ejemplo, cuándo se van a cometer crímenes incluso antes de que se cometan, como ocurría en la película `Minority Report' \cite{movie:minorityreport}, pero ahora ya es una realidad. Otro ejemplo es el reciente uso de \textit{big data} en China como un sistema disciplinario \cite{art:chinabigdata}. En ambos ejemplos destaca la importancia de los datos demográficos y sociológicos que hasta ahora a penas se han tratado en el ámbito español.

Un tema que está siendo investigado actualmente es comprender las causas por las que se producen crímenes. Por ello, se han realizado estudios para investigar las tasas de criminalidad utilizando técnicas de minería de datos \cite{art:crimeprediction} \cite{art:crimeprediction2} a partir de datos demográficos y sociológicos con el objetivo de determinar los lugares donde asignar más recursos para reducir estos índices de criminalidad.

Se dice que los datos se están convirtiendo en el nuevo petróleo, en el sentido de que con unos datos de buena calidad pueden descubrirse relaciones insospechadas hasta el momento. Una fuente importante de datos son los datos públicos, pero tiene inconvenientes, estos datos se obtienen de múltiples fuentes, por lo que cada una de ellas nos proporciona la información organizada de forma dirente y en distintos formatos, también podrían ser accesibles de modo diferente, lo que nos complica la laboral al analizarlos.

El objetivo de este trabajo es llevar estos estudios al ámbito español y unificar estas fuentes de datos de forma que el acceso sea uniforme, mucho más sencillo y permitiendo combinar fuentes de datos de forma directa. Se ha realizado integrando varias fuentes de datos públicas utilizando técnicas de web scraping (\ref{webscraping}). Para  después visualizar estos datos de manera sencilla, permitir crear mapas temáticos (\ref{mapascoropleticos}) y facilitar la exportación de los datos para después ser utilizados en procesos de minería de datos de forma directa.

Por último, agradecimiento a la idea original de Roberto Cuesta, capitán de la guardia civil y estudiante del doctorado de la Universidad de Burgos, con una tesis enfocada a la aplicación de técnicas de minería de datos en el ámbito criminológico.

\subsection{Fuentes de datos}

A continuación se van a listar las fuentes de datos de organismos públicos que se han integrado.

\begin{itemize}
	\tightlist
	\item
	Servicio Público de Empleo Estatal (SEPE)
	\begin{itemize}
		\tightlist
		\item
		Paro registrado por municipio\footnote{\href{https://datos.gob.es/catalogo/e00142804-paro-registrado-por-municipios}{https://datos.gob.es/catalogo/e00142804-paro-registrado-por-municipios}}.
		\item
		Contratos registrados por municipio\footnote{\href{http://datos.gob.es/es/catalogo/e00142804-contratos-por-municipios}{http://datos.gob.es/es/catalogo/e00142804-contratos-por-municipios}}.
		\item
		Demandantes de empleo por municipio\footnote{\href{http://datos.gob.es/es/catalogo/e00142804-demandantes-de-empleo-por-municipios}{http://datos.gob.es/es/catalogo/e00142804-demandantes-de-empleo-por-municipios}}.
	\end{itemize}
	\item 
	Estadísticas de la renta por municipio (Agencia tributaria)\footnote{\href{https://www.agenciatributaria.es/AEAT.internet/datosabiertos/catalogo/hacienda/Estadistica\_de\_los\_declarantes\_del\_IRPF\_por\_municipios.shtml}{https://www.agenciatributaria.es/AEAT.internet/datosabiertos/catalogo/hacienda/ Estadistica\_de\_los\_declarantes\_del\_IRPF\_por\_municipios.shtml}}.
	\item
	Instituto Nacional de Estadística (INE)
	\begin{itemize}
		\item
		Estadísticas de población por sexo y edad\footnote{\href{http://www.ine.es/jaxi/Tabla.htm?path=/t20/e245/p05/a2011/l0/\&file=00000001.px\&L=0}{http://www.ine.es/jaxi/Tabla.htm?path=/t20/e245/p05/a2011/l0/\&file=00000001.px\&L=0}}.
		\item
		Relación de municipios y códigos por provincias\footnote{\href{http://www.ine.es/daco/daco42/codmun/codmunmapa.htm}{http://www.ine.es/daco/daco42/codmun/codmunmapa.htm}}.
	\end{itemize}
	\item
	Ministerio del interior (MIR)
	\begin{itemize}
		\item 
		Resultados municipales en las elecciones al congreso\footnote{\href{http://www.infoelectoral.mir.es/infoelectoral/min/areaDescarga.html?method=inicio}{http://www.infoelectoral.mir.es/infoelectoral/min/areaDescarga.html?method=inicio}}.
	\end{itemize}
\end{itemize}


\subsection{Material entregado}

Material adjunto a la memoria:

\begin{itemize}
	\item 
	Aplicación web en Flask.
	\item 
	Scripts para la integración de fuentes de datos.
	\item
	Scripts para la ejecución.
	\item 
	Aplicación de escritorio para juntar csv.
	\item 
	Anexos.
\end{itemize}

Recursos disponibles en internet:

\begin{itemize}
	\item 
	\foothref{Página web del proyecto}{https://tfg-datos-publicos.nanoapp.io/}.
	\item
	\foothref{Repositorio del proyecto}{https://github.com/IvanArjona/TFG-Datos-publicos}.
\end{itemize}

\capitulo{2}{Objetivos del proyecto}

Este apartado explica de forma precisa y concisa cuales son los objetivos que se persiguen con la realización del proyecto. Se puede distinguir entre los objetivos marcados por los requisitos del software a construir y los objetivos de carácter técnico que plantea a la hora de llevar a la práctica el proyecto.

\capitulo{3}{Conceptos teóricos}

En aquellos proyectos que necesiten para su comprensión y desarrollo de unos conceptos teóricos de una determinada materia o de un determinado dominio de conocimiento, debe existir un apartado que sintetice dichos conceptos.


\section{Bases de datos no relacionales}

Las bases de datos no relacionales \cite{wiki:nosql}, también llamadas NoSQL (`Not Only SQL') son bases de datos optimiazadas para ser utilizadas con modelos de datos sin esquema y potencialmente escalables.

A diferencia de las bases de datos relacionales, aquí no hay tables, esquemas ni relaciones, sino que los datos pueden almacenarse con cualquier esquema sin tener que seguir todos la misma estructura.

En este tipo de bases de datos tenemos colecciones en lugar de tablas y dentro de estas colecciones tenemos documentos. Estos documentos pueden contener cualquier información y se almacenan siguiendo un esquema json.

Una de las razones importantes por las que usar bases de datos NoSQL en lugar de relacionales es la alta velocidad de consulta y la gran capacidad de escalabilidad.


\section{Web scraping} \label{webscraping}

Web Scraping \cite{wiki:webscraping} es una técnica para extraer información de sitios web directamente de su código fuente sin utilizar APIs\footnote{Interfaz de programación de aplicaciones} proporcionadas por el propio sitio.

Lo que hacemos con un \textit{web scraper} es buscar información dentro de un documento web siguiendo ciertos patrones en la estructura de su código fuente y extraer esta información a nuestro entorno local.

Las técnicas de web scraping se centran en transformar datos sin estructura de una página web en datos estructurados para poder ser almacenados y analizados posteriormente. Por ejemplo, en una base de datos, hojas de cálculo de dataframes de pandas.

\section{Mapas coropléticos} \label{mapascoropleticos}

Un mapa coroplético \cite{wiki:mapascoropleticos} es un mapa topológico dividido en regiones en el que cada una de estas regiones se pinta de un color de acuerdo a una medida estadística.

Para ilustrarlo mejor pondremos el ejemplo del siguiente mapa en la figura \ref{fig:conceptos/mapa-paro}. En él se compara el paro de todas las provincias de España pintando con colores más cálidos las provincias con mayor porcentaje de paro y con colores más fríos las provincias con menor porcentaje de paro.

\imagencite{conceptos/mapa-paro}{Mapa coroplético del paro en las provincias españolas}{art:mapaparo}


\capitulo{4}{Técnicas y herramientas}

En esta sección se van a explicar, por un lado, las técnicas que se han segueido durante el desarrollo del proyecto y por otro, las herramientas utilizadas, que se han dividido en herramientas de desarrollo, gestión y documentación.

\section{Técnicas}

\subsection{Metodología ágil}

Este proyecto se ha realizado siguiendo los principios del manifiesto ágil~\cite{art:agilemanifesto}. En concreto se ha seguido la metodología \textit{scrum} \cite{book:scrum} con sprints normalmente de una semana y reuniones semanales de planificación y revisión.

Al utilizar esta metodología se persigue utilizar una estrategia de desarrollo incremental, en la que al final de cada sprint se tiene como resultado un incremento del software entregable. También se permite que los requisitos cambien a lo largo del proyecto, en lugar de fijarse al principio, consiguiendo un software de mejor calidad.

Un software creado con unos requisitos claros y fijos desde el principio puede ser un software de muy alta calidad y podría realizarse con una metodología tradicional, pero para proyectos exploratorios como este, en los que los requisitos pueden cambiar, es mejor utilizar una metodología ágil.

Para seguir esta metodología se han utilizado las \foothref{issues de Github}{https://github.com/IvanArjona/TFG-Datos-publicos/issues} (\ref{github}) y los tableros Kanban de \foothref{ZenHub}{https://www.zenhub.com/} (\ref{zenhub}) para organizar la pila de tareas a realizar en cada sprint.

\subsection{DevOps}

DevOps \cite{misc:devops} es una metodología para la creación de software en la que se integra el desarrollo de software y la administración de sistemas.

Se ha elegido esta metodología para conseguir desplegar rápidamente todos los cambios en el servidor web. Apoyándonos en herramientas como Nanobox (\ref{nanobox}) para facilitar este proceso.

\section{Herramientas}

\subsection{Herramientas de desarrollo}

A continuación se van a explicar las herramientas utilizadas para desarrollar el proyecto.

\subsubsection{MongoDB}

\begin{itemize}
	\tightlist
	\item
	Herramientas consideradas:
	\href{http://basho.com/riak/}{Riak}, 
	\href{http://cassandra.apache.org/}{Cassandra},
	\href{https://www.mongodb.com/}{MongoDB}, 
	\href{http://leveldb.org/}{LevelDB}.
	\item
	Herramienta elegida:
	\href{https://www.mongodb.com/}{MongoDB}.
\end{itemize}

\foothref{MongoDB}{https://www.mongodb.com/} \cite{misc:mongodb} es un sistema de bases de datos NoSQL. En esta herramienta los datos se guardan en forma de documentos con un esquema similar a JSON. Con este sistema se consigue una consulta de datos más rápida. Además se ha usado PyMongo \cite{docs:pymongo} como herramienta para integrar MongoDB en Python.

Tanto Riak como Cassandra son también bases de datos NoSQL, se descartaron porque no ofrecen soporte para equipos con sistema operativo Windows y no son bases de datos de documentos, que se cree es lo más conveniente para este proyecto. 

LevelDB es una base de datos NoSQL de pares clave-valor, al igual que Riak, esta herramienta se descartó porque no se cree conveniente utilizar pares clave-valor para un proyecto como este y no hay tantos ejemplos en la documentación como en las otras herramientas.

Podría haberse utilizado bases de datos Relacionales, pero la razón por la que se eligió MongoDB es por la facilidad de almacenar los dataframes de Pandas y, sobre todo, la motivación principal fue la inquietud por formarme en esta tecnología por lo extendida que está y porque no se ve nada en la carrera.

\subsubsection{DigitalOcean} \label{digitalocean}

\begin{itemize}
	\tightlist
	\item
	Herramientas consideradas:
	\href{https://www.heroku.com/}{Heroku}, 
	\href{https://www.pythonanywhere.com/}{PythonAnywhere},
	\href{https://www.digitalocean.com/}{DigitalOcean}, 
	\href{https://aws.amazon.com/es/}{Amazon Web Services}.
	\item
	Herramienta elegida:
	\href{https://www.digitalocean.com/}{DigitalOcean}.
\end{itemize}

\foothref{DigitalOcean}{https://www.digitalocean.com/} es un proveedor de servidores privados, por ello podemos hacer lo que queramos con el servidor sin limitaciones más allá de la capacidad de procesamiento y memoria RAM. Se ha elegido este servicio porque nos da total libertad y se puede probar gratuitamente con \foothref{GitHub Education}{https://education.github.com/}~\cite{docs:digitalocean}.

Una alternativa que se consideró y de hecho, se probó es Heroku, en este caso se instala el entorno necesario de forma automática. El problema es que la base de datos en la capa gratuita sólo puede pesar 500MB como máximo y no es suficiente para este proyecto.

PythonAnywhere es un hosting para aplicaciones web en Python, el problema con este proveedor es que no ofrece bases de datos locales, por lo que habría que utilizar una remota. La única gratuita que se ha encontrado es \foothref{mLab}{https://mlab.com/}, la misma que usa Heroku, por lo que volvemos al mismo problema del límite de tamaño.

Por último, se consideró utilizar una instancia de \foothref{Amazon AWS EC2}{https://aws.amazon.com/es/ec2/}. Es muy similar a DigitalOcean, se eligió el primero porque es más sencillo de utilizar.

\subsubsection{Nanobox} \label{nanobox}
\begin{itemize}
	\tightlist
	\item
	Herramientas consideradas:
	\href{https://www.heroku.com/}{Heroku}, 
	\href{https://nanobox.io/}{Nanobox}.
	\item
	Herramienta elegida:
	\href{https://nanobox.io/}{Nanobox}.
\end{itemize}

\foothref{Nanobox}{https://nanobox.io/} es una herramienta que nos permite desplegar nuestra aplicación sin centrarnos en la infraestructura del servidor \cite{docs:nanobox}.

Para ello enlazamos nuestra cuenta de un proveedor en la nube (en este caso DigitalOcean) y nanobox se encargará de instalar el sistema operativo, configurarlo, instalar nuestra aplicación y sus dependencias y ejecutarla.

Heroku es un servicio muy similar, con la diferencia de que no podemos utilizar servidores en la nube externos, se descartó por lo ya explicado en el punto anterior.

\subsubsection{Flask}

\begin{itemize}
	\tightlist
	\item
	Herramientas consideradas:
	\href{http://flask.pocoo.org/}{Flask}, 
	\href{https://www.djangoproject.com/}{Django}.
	\item
	Herramienta elegida:
	\href{http://flask.pocoo.org/}{Flask}.
\end{itemize}

Como uno de los objetivos del proyecto es el de crear una página web, se han considerado varios frameworks web para Python. Entre ellos se ha seleccionado Flask.

\foothref{Flask}{http://flask.pocoo.org/} es un framework fácil de utilizar y muy flexible. No nos fuerza a utilizar una metodología específica y podemos organizar la aplicación con la estructura que queramos (a diferencia de \textit{Django}) \cite{book:flask}.

Además se incluyen herramientas para desplegar el servidor de desarrollo, para realizar pruebas de la aplicación y para hacer \textit{APIs REST}.

\subsubsection{Folium} \label{folium}

\begin{itemize}
	\tightlist
	\item
	Herramientas consideradas:
	\href{https://plot.ly/python/maps/}{Plotly maps}, 
	\href{http://python-visualization.github.io/folium/}{Folium},
	\href{https://leafletjs.com/}{Leaflet}.
	\item
	Herramienta elegida:
	\href{http://python-visualization.github.io/folium/}{Folium}.
\end{itemize}

Un aspecto importante de este proyecto es la representación de datos en el mapa.

Para ello se han considerado varias herramientas, de las cuales se ha elegido \foothref{Folium}{http://python-visualization.github.io/folium/} \cite{docs:folium}. Este paquete nos permite visualizar datos manipulados con Python y visualizarlos como mapas utilizando para ello \textit{LeafletJS} \cite{docs:leaflet}.

Se ha elegido esta herramienta porque nos permite representar mapas coropléticos (\ref{mapascoropleticos}) y mapas con clusters de datos \cite{misc:foliumcluster} utilizando estructuras de datos geográficas personalizadas (\textit{GeoJSON}) \cite{docs:geojson}.

\subsubsection{Dynatable}

\begin{itemize}
	\tightlist
	\item
	Herramientas consideradas:
	\href{https://www.dynatable.com/}{Dynatable}, 
	\href{https://datatables.net/}{Datatables}.
	\item
	Herramienta elegida:
	\href{https://www.dynatable.com/}{Dynatable}.
\end{itemize}

\foothref{Dynatable}{https://www.dynatable.com/} es un framework de Javascript para visualizar tablas de una manera más clara y ordenada \cite{docs:dynatable}.

Este framework nos permite paginar los resultados, ordenar por alguna columna y buscar por cualquier campo dentro de la tabla. Todo esto sin tener que recargar la página.

\subsubsection{Bootstrap}

\begin{itemize}
	\tightlist
	\item
	Herramientas consideradas:
	\href{https://getbootstrap.com/}{Bootstrap}, 
	\href{https://foundation.zurb.com/}{Foundation}.
	\item
	Herramienta elegida:
	\href{https://getbootstrap.com/}{Bootstrap}.
\end{itemize}

\foothref{Bootstrap}{https://getbootstrap.com/} es un framework para desarrollar páginas web con HTML, CSS y Javascript para facilitar el diseño de páginas responsive, que se adapten al tamaño de la pantalla del dispositivo.

Bootstrap nos ofrece una serie de clases de CSS predefinidas por lo que podemos crear la web del proyecto prácticamente sin tocar hojas de estilos CSS. Otra característica es que dispone de funciones de Javascript para llamar a eventos de ciertos elementos, como puede ser al pulsar un item de un menú desplegable.

\subsubsection{PyCharm}

\begin{itemize}
	\tightlist
	\item
	Herramientas consideradas:
	\href{https://code.visualstudio.com/}{Visual Studio Code}, 
	\href{https://www.jetbrains.com/pycharm/}{PyCharm}.
	\item
	Herramienta elegida:
	\href{https://www.jetbrains.com/pycharm/}{PyCharm}.
\end{itemize}

\foothref{PyCharm}{https://www.jetbrains.com/pycharm/} es un IDE (\textit{Entorno de Desarrollo Integrado}) para Python basado en IntelliJ. Posee herramientas para ayudarnos con el desarrollo como el autocompletado y herramientas para refactorizar de forma automática.

Se ha considerado porque, a diferencia de otros IDEs similares, tiene soporte para Flask y MongoDB, lo que nos facilita programar más rápidamente.

En este proyecto se ha utilizado la versión \textit{Professional}, pero con la versión \textit{Community} gratuita sería suficiente.

\subsection{Herramientas de gestión}

A continuación se van a describir las herramientas que se han utilizado para gestionar el proyecto.

\subsubsection{Github} \label{github}

\begin{itemize}
	\tightlist
	\item
	Herramientas consideradas:
	\href{https://github.com/}{Github}, 
	\href{https://bitbucket.org/product}{BitBucket},
	\href{https://about.gitlab.com/}{GitLab}.
	\item
	Herramienta elegida:
	\href{https://github.com/}{Github}.
\end{itemize}

\foothref{GitHub}{https://github.com/} es un servicio web para alojar proyectos utilizando un sistema de control de versiones \textit{Git}.

Se ha utilizado esta herramienta para alojar el \href{https://github.com/IvanArjona/TFG-Datos-publicos}{repositorio del proyecto} llevando un registro de todos los cambios realizados desde el inicio.

Además se han utilizado las \textit{Issues} y \textit{Milestones} de GitHub para planificar las tareas a realizar en cada sprint.

\subsubsection{ZenHub} \label{zenhub}

\begin{itemize}
	\tightlist
	\item
	Herramientas consideradas:
	\href{https://help.github.com/articles/about-project-boards/}{Github Projects}, 
	\href{https://www.zenhub.com/}{ZenHub},
	\href{https://www.gitkraken.com/glo}{GitKraken Glo}.
	\item
	Herramienta elegida:
	\href{https://www.zenhub.com/}{ZenHub}.
\end{itemize}

\foothref{ZenHub}{https://www.zenhub.com/} es una herramienta de gestión de proyectos integrada en Github mediante una extensión para el navegador. Nos proporciona un tablero canvas sobre el que mover las \textit{issues} entre varios estados.

Las tareas se pueden estimar estableciendo un número de `\textit{Story points}', basados en la serie de Fibonacci en función del trabajo que se cree que se va a necesitar. Esto ha sido especialmente útil para planificar el tiempo empleado en cada tarea \cite{misc:zenhubestimate}.

También se ha utilizado esta herramienta para extraer los \textit{gráficos burndown} de cada sprint presentes en el anexo.

\subsubsection{GitKraken}

\begin{itemize}
	\tightlist
	\item
	Herramientas consideradas:
	\href{https://www.gitkraken.com/}{GitKraken}, 
	\href{https://desktop.github.com/}{Github Desktop},
	\href{https://www.sourcetreeapp.com/}{Sourcetree}.
	\item
	Herramienta elegida:
	\href{https://www.gitkraken.com/}{GitKraken}.
\end{itemize}

\foothref{GitKraken}{https://www.gitkraken.com/} es un cliente con interfaz gráfica para gestión de proyectos \textit{Git} para Windows, Mac y Linux. Nos permite realizar todas las acciones posibles con git, pero de manera mucho más sencilla y visual.

Permite sincronizarse con Github y otros servicios de alojamiento de repositorios de forma automática, de forma que el contenido del repositorio local y el remoto sea consistente.

\subsubsection{Codacy}

\begin{itemize}
	\tightlist
	\item
	Herramientas consideradas:
	\href{https://www.sonarqube.org/}{SonarQube}, 
	\href{https://codeclimate.com/}{Code Climate},
	\href{https://www.codacy.com/}{Codacy}.
	\item
	Herramienta elegida:
	\href{https://www.codacy.com/}{Codacy}.
\end{itemize}

\foothref{Codacy}{https://www.codacy.com/} es una herramienta de análisis de calidad de código que ayuda a los desarrolladores a realizar código de más rápido y de mayor calidad.

Tiene métricas para medir la complejidad ciclomática, código duplicado, cubrimiento por tests y estadísticas para cada commit o pull request \cite{misc:codacygithub}.

Durante el desarrollo de este proyecto se ha intentado mantener siempre una Certificación A en Codacy con todas las métricas al 100\% para asegurar la calidad del código.

\imagencite{herramientas/codacy}{Certificado de Codacy del proyecto}{misc:codacycertificado}

\subsection {Herramientas de documentación}

\subsubsection{LaTeX}

\begin{itemize}
	\tightlist
	\item
	Herramientas consideradas:
	\href{https://products.office.com/es-es/word}{Microsoft Word}, 
	\href{https://www.latex-project.org/}{LaTeX},
	\href{https://es.libreoffice.org/}{LibreOffice},
	\href{https://www.openoffice.org/es/}{OpenOffice}.
	\item
	Herramienta elegida:
	\foothref{LaTeX}{https://www.latex-project.org/}.
\end{itemize}

Esta memoria está escrita en \LaTeX. Se trata de un sistema de preparación de documentos para la creación de documentos con una alta calidad tipográfica. Se usa frecuentemente en la escritura de libros y artículos científicos, entre otra cosas por la facilidad de incluir expresiones matemáticas \cite{wiki:latex} \cite{wiki:latexwikibooks}.

Para utilizar esta herramienta se ha utilizado \foothref{TeXstudio}{https://www.texstudio.org/}. Es un editor gratuito y multiplataforma para editar documentos escritos en Latex. Incluye iconos para localizar los comandos más comunes, corrector ortográfico, visor de PDF y descarga de paquetes de forma automática.

Este documento se ha escrito utilizando la siguiente plantilla de Latex~\cite{misc:plantillalatex}.

\subsubsection{StarUML}

\begin{itemize}
	\tightlist
	\item
	Herramientas consideradas:
	\href{http://staruml.io/}{StarUML}, 
	\href{http://argouml.tigris.org/}{ArgoUML},
	\href{http://astah.net/}{Astah}.
	\item
	Herramienta elegida:
	\href{http://staruml.io/}{StarUML}.
\end{itemize}

\foothref{StarUML}{http://staruml.io/} es una herramienta para modelar diagramas UML que nos permite crear diagramas de clases, objetos, casos de uso, componente, estructurales, de comunicación, de estado y de actividad y entidad-relación.

Funciona de manera gráfica, arrastrando objetos y conectándolos entre ellos permitiendo añadir atributos tanto a los objetos como a las relaciones.

\capitulo{5}{Aspectos relevantes del desarrollo del proyecto}

Este apartado pretende recoger los aspectos más interesantes del desarrollo del proyecto, comentados por los autores del mismo.
Debe incluir desde la exposición del ciclo de vida utilizado, hasta los detalles de mayor relevancia de las fases de análisis, diseño e implementación.
Se busca que no sea una mera operación de copiar y pegar diagramas y extractos del código fuente, sino que realmente se justifiquen los caminos de solución que se han tomado, especialmente aquellos que no sean triviales.
Puede ser el lugar más adecuado para documentar los aspectos más interesantes del diseño y de la implementación, con un mayor hincapié en aspectos tales como el tipo de arquitectura elegido, los índices de las tablas de la base de datos, normalización y desnormalización, distribución en ficheros3, reglas de negocio dentro de las bases de datos (EDVHV GH GDWRV DFWLYDV), aspectos de desarrollo relacionados con el WWW...
Este apartado, debe convertirse en el resumen de la experiencia práctica del proyecto, y por sí mismo justifica que la memoria se convierta en un documento útil, fuente de referencia para los autores, los tutores y futuros alumnos.

\capitulo{6}{Trabajos relacionados}

%TODO: papers

\section{Mapa del paro}

\foothref{Mapa del paro}{http://mapadelparo.com/} es una página web que nos permite visualizar el paro por municipios, provincias y comunidades en un mapa coroplético. Ver figura~\ref{fig:relacionados/mapadelparo}.

Los datos del paro contenidos en esta página web también están en este proyecto, con la ventaja de que además se pueden combinar con otros datos públicos y se puede descargar la información.

\imagen{relacionados/mapadelparo}{Paro por comunidades autónomas}

\section{Ciudatos}

\foothref{Ciudatos}{http://www.ciudatos.com/visualizacion} nos permite visualizar varios datos públicos de Colombia mediante gráficas de barras y posicionarlos en el mapa. Ver figura \ref{fig:relacionados/ciudatos}.

Esta herramienta no nos permite visualizar los datos con mapas coropléticos, simplemente sitúa los gráficos de barras de cada región en su correspondiente sitio.

Como ventaja nuestra aplicación tiene datos de España y permite visualizar varios datos diferentes entre ellos al mismo tiempo.

\imagen{relacionados/ciudatos}{Bienestar por ciudades (Colombia)}

\section{Eurostat}

Otra representación de datos públicos en mapas es \foothref{Eurostat}{http://ec.europa.eu/eurostat/en/web/lfs/statistics-illustrated}.

En esta representación se muestra en el mapa los porcentajes de empleo en cada país de Europa. Ver figura \ref{fig:relacionados/eurostat}.

Estos datos se escapan del alcance al que se quería llegar en este trabajo, pero podría considerarse para integrarlos en algún proyecto futuro.

\imagen{relacionados/eurostat}{Porcentaje de empleo por país (Europa)}

\capitulo{7}{Conclusiones y Líneas de trabajo futuras}

En este capítulo se exponen las conclusiones derivadas del desarrollo del proyecto y comentarios para continuar con el trabajo.

\section{Conclusiones}

\section{Líneas de trabajo futuras}

Este trabajo es el inicio de un proyecto más completo y más ambicioso que se pretende continuar en trabajos de fin de grado futuros. A continuación se muestra una lista de posibles tareas a realizar en trabajos futuros:

\begin{itemize}
	\item Mejoras las columnas calculadas:
	\begin{itemize}
		\item Validar si la consulta es correcta dinámicamente con \textit{Javascript}.
		\item Permitir consultar más de una columna calculada.
		\item Permitir elegir el nombre de las columnas calculadas.
		\item Mejorar el autocompletado al introducir operadores.
	\end{itemize}
	\item Hacer que los datos estén disponibles en la web semántica siguiendo la descripción de RDF~\cite{misc:informationserviceengineering}.
	\item Mejorar la visualización de mapas: mostrar el valor de los atributos al pasar el ratón sobre una región, hacer que funcione el renderizado de mapas en Chrome. Puede que haya que cambiar de framework para representar mapas. 
	\item Integrar más fuentes de datos estatales. Como por ejemplo:
	\begin{itemize}
		\item Matrimonios por provincia de residencia del matrimonio, mes de celebración y forma de celebración del matrimonio
		\footnote{\href{http://www.ine.es/dynt3/inebase/index.htm?padre=3406}{http://www.ine.es/dynt3/inebase/index.htm?padre=3406}}.
		\item Balances de criminalidad
		\footnote{\href{http://www.interior.gob.es/prensa/balances-e-informes/2017}{http://www.interior.gob.es/prensa/balances-e-informes/2017}}.
		\item Estadísticas catastrales
		\footnote{\href{http://www.catastro.minhap.es/esp/estadisticas\_2.asp}{http://www.catastro.minhap.es/esp/estadisticas\_2.asp}}.
		\item Información meteorológica
		\footnote{\href{http://www.aemet.es/es/datos\_abiertos/catalogo}{http://www.aemet.es/es/datos\_abiertos/catalogo}}.
	\end{itemize}
	\item Integrar fuentes de datos de empresas privadas:
	\begin{itemize}
		\item Tendencias de búsqueda como \textit{Google Trends}
		\footnote{\href{https://trends.google.es/trends/?geo=ES}{https://trends.google.es/trends/?geo=ES}}.
		\item Porcentajes de audiencia por municipio (televisión, radio, periódico).
	\end{itemize}
\end{itemize}


\bibliographystyle{plain}
\bibliography{bibliografia}

\end{document}
